\subsubsection{IP Interrupt Enable Register (IPIER)}

Address offset: \textbf{0x08}
\newline
Reset value: 0x0000 0000

\begin{figure}[H]
    \centering
    \begin{bytefield}[bitwidth=\bflenhalfword, endianness=big]{16}
        \bitheader[lsb=16]{16-31} \\
        \bitbox{16}{Reserved}
    \end{bytefield}
\end{figure}

\begin{figure}[H]
    \centering
    \begin{bytefield}[bitwidth=\bflenhalfword, endianness=big]{16}
        \bitheader{0-15} \\
        \bitbox{15}{Reserved} &
        \bitbox{1}{FIE} \\
        \bitbox[]{15}{}
        \bitbox[]{1}{R/W}
    \end{bytefield}
\end{figure}

\begin{itemize}
    \item \textbf{Bit 31:1 - Reserved Bits}
    \item \textbf{Bit 0 - FIE: Frame Data Empty Interrupt Enable}\\
    Wenn dieses Bit auf 1 gesetzt wird, ist der Frame Data Empty Interrupt aktiviert.
    Er wird ausgelöst, wenn das FDE Bit in IPISR gesetzt wird.
    Zusätzlich müssen global die Interrupts aktiviert sein indem das GIE Bit im GIER Register gesetzt ist.
\end{itemize}